\documentclass[12pt]{article}
\usepackage[a4paper,margin=2cm]{geometry}
\usepackage{amsmath,amssymb}
\usepackage{xepersian}
\settextfont[
  Path=../1401-2/HW2 - official solution/,
  Extension=.ttf,
  BoldFont=*Bd,
  ItalicFont=*It,
  BoldItalicFont=*BdIt
]{XB Niloofar}
\setdigitfont[
  Path=../1401-2/HW2 - official solution/,
  Extension=.ttf,
  BoldFont=*Bd,
  ItalicFont=*It,
  BoldItalicFont=*BdIt
]{XB Niloofar}
\setlength{\parindent}{0pt}
\setlength{\parskip}{0.55em}

\begin{document}
\begin{center}
  {\Large\textbf{اصلاحیهٔ پاسخ رسمی تمرین ششم ـ نیم‌سال ۱۴۰۲--۲}}
\end{center}

این برگه فقط موارد نادرست پاسخ منتشرشده را اصلاح می‌کند و باید همراه فایل اصلی خوانده شود.

\section*{سؤال ۲}
معکوس شبهِ مور--پنروز هر ماتریس وجود دارد و \textbf{یکتا} است. آنچه ممکن است یکتا نباشد، بردار ضرایب \(\alpha\) در نمایش
\[
  w^*=X^{T}\alpha
\]
یا یک معکوس تعمیم‌یافتهٔ دلخواه است. بنابراین استدلال پاسخ اصلی که «معکوس شبه یکتا نیست» را به‌کار می‌برد باید با این تمایز اصلاح شود.

\section*{سؤال ۳}
با قرارداد پاسخ اصلی که
\(\theta^{(k)}=\Phi^T\beta^{(k)}\)
است، جایگذاری در گام گرادیان نزولی می‌دهد
\[
 \theta^{(k+1)}
 =\Phi^T\!\left(
 \beta^{(k)}
 -2\alpha\,\Phi\Phi^T\beta^{(k)}
 +2\alpha y
 \right).
\]
پس رابطهٔ درست ضرایب برابر است با
\[
  \boxed{\;
  \beta^{(k+1)}
  =\beta^{(k)}
  -2\alpha\,\Phi\Phi^T\beta^{(k)}
  +2\alpha y
  \;}
\]
یا، با \(K=\Phi\Phi^T\)،
\(\beta^{(k+1)}=\beta^{(k)}-2\alpha K\beta^{(k)}+2\alpha y\).
ترتیب \(\Phi^T\Phi\) در پاسخ اصلی از نظر ابعادی نادرست است.

\section*{سؤال ۵}
غلبهٔ قطری و مثبت‌معین‌بودن متقارن، شرط‌های \textbf{کافی} برای همگرایی گاوس--سایدل‌اند، نه شرط‌های لازم. همچنین ادعای پاسخ اصلی دربارهٔ سطر اول نادرست است، زیرا
\[
  25>5+1.
\]
البته غلبه‌ی قطری در سطرهای دوم و سوم برقرار نیست، زیرا
\[
  8<64+1,
  \qquad
  1<144+12.
\]
برای ماتریس
\[
 A=\begin{bmatrix}
 25&5&1\\
 64&8&1\\
 144&12&1
 \end{bmatrix}=L+U
\]
معیار قطعی، شعاع طیفی ماتریس تکرار
\(B=-L^{-1}U\)
است. مقادیر ویژهٔ آن تقریباً
\[
  0,\qquad 0.8668019,\qquad 4.1531981
\]
هستند؛ بنابراین
\[
  \rho(B)\simeq4.1531981>1
\]
و روش برای این ترتیب معادلات واگراست. این محاسبه، نه صرفاً نبودن یک شرط کافی، دلیل درست واگرایی است.
\end{document}
